\section{The Loss of God}

The chief punishment of hell is the loss of God. The expression can seem brief beside images of fire, but it names the deepest deprivation a creature can suffer. God is the source of every good and the end for which we were made. The blessed behold him and live in his friendship, with every holy desire fulfilled in him. To be excluded forever from that communion is to lose the good to which the whole person is ordered. Nothing left within the created order can make up for it.

Separation from God does not mean that the condemned become independent of their Creator. They remain creatures sustained by his power. What is lost is the blessed possession of him, the communion of love and glory. This makes the loss more searching than exile from a beloved earthly place. A traveler can find another home; a person deprived of an earthly good may discover consolation in another. No other final end waits beside God to receive the soul and make it blessed.

St.~John Chrysostom develops this truth in his twenty-third homily on Matthew. The homily has followed Christ's teaching about the narrow gate, false prophets, and the tree that is cut down and cast into the fire. Chrysostom observes that the condemned person suffers both the burning and exclusion from the kingdom. He calls the latter the more grievous punishment. He also recognizes why his hearers find it difficult to grasp: they do not adequately know the greatness of the blessedness they stand to lose (\emph{Homilies on Matthew} 23.9).

That difficulty is itself a subject for meditation. We readily understand the loss of health, money, work, companionship, or honor, because we have tasted those goods and organized our lives around them. The loss of God can seem distant if prayer is distant, if the Eucharist is received without attention, or if heaven is seldom considered. Our weak apprehension does not measure the good's smallness. It reveals how much the heart still needs to learn to desire its true end.

Chrysostom makes the deprivation personal by placing Christ before his hearers. He asks them to think of a father who would bear any loss rather than be parted from a beloved child. Even that strong attachment cannot equal the good of being with Christ. The damned lose the blessed company of the Savior. Chrysostom dwells on the face of the Lord, on hearing oneself rejected by him, and on the love with which that Lord had first sought the sinner.

Chrysostom recalls how Christ has created and sustained us, suffered for us, given his Spirit, and offered his kingdom. These gifts establish the sinner's debt of gratitude. The Lord who pronounces the sentence has already sought the person he judges; the face lost forever is the face of his benefactor. A Christian considering that deprivation also remembers his own baptism and forgiveness. Repeated invitations to return make continued resistance a refusal of love already received.

At this point Chrysostom turns to the hungry Christ. The Savior has given himself for the sinner, and the sinner will not give him bread. The Lord's identification with the poor in Matthew 25 brings eternal loss into the present relation between giver and beggar. The hearer cannot safely admire the doctrine and remain unconcerned about the hungry person. The Christ whose glorious face we hope to see has entrusted a present claim upon our charity to the needy (23.9).

Chrysostom then compares the pursuit of wealth with the relief of suffering. One man spends his efforts acquiring fields and possessions; another uses his goods to relieve need, heal sickness, and release people from distress. The second becomes a refuge to others. He compares worldly power with patience under injury as well. The powerful person may be swollen with pride, while the person who bears an insult meekly possesses a dignity that death cannot strip away (23.10--12). Wealth and influence can command admiration now, but the use of them reveals whether their possessor is preparing for the kingdom.

This is a demanding account of charity because it concerns what we love. We may prefer to be known as benefactors while avoiding the person whose need is inconvenient. We may desire the respect accorded to virtue while refusing the humiliation that would let a quarrel end. Chrysostom's treatment of the loss of Christ presses beneath those arrangements. The heart must be prepared for the Lord himself. Possessions, influence, and public esteem are to serve that end, and become dangerous when we protect them at its expense.

St.~Alphonsus begins his consideration of hell with the two movements of sin: turning away from God and turning toward a creature. The created good is sought against the order established by its Creator. Accordingly he distinguishes the pain of loss from the other punishments, placing the loss of God first in gravity (\emph{Preparation for Death}, XXVI.1). His distinction explains what the sinner is choosing. We seek a particular satisfaction, while the grace and friendship we surrender can receive scarcely a moment's serious attention.

In the third point, Alphonsus considers how grief over a loss grows with our estimate of what has been lost. His example is a precious jewel: its value makes its loss more keenly felt. But God exceeds every scale of exchange. The soul's final deprivation cannot be understood as losing one valuable object while retaining enough others to be content. The good lost is infinite. Alphonsus brings the consideration back to Christ's sentence of departure and to the terrible permanence of never seeing him in glory (XXVI.3).

The saint then turns his meditation toward the grace still available. His prayers recognize that the living sinner can yet love God, repent, and recover friendship with him. The possibility of loving him now is already a great mercy. We can share that gratitude without pretending that we have measured the depths of eternal loss. To recognize God as our supreme good is enough to make a beginning: the attachment which costs his friendship must be renounced, and the love he commands must be sought from him.

We can dwell upon this loss without continually seeking a more alarming image. Consider the good of belonging to Christ and the attachment for which we endanger that friendship. The attraction of the attachment is familiar; we have rehearsed its advantages whenever temptation returns. Prayer gives time to consider what we have left out of that calculation. A satisfaction which passes cannot compensate for the everlasting deprivation of God. The choice becomes clearer when the will considers both what it wants immediately and what it risks losing.

